\section{Introduction}\label{sec:intro}
Power systems generate interdependent streams of forecasts, prices, operating quantities and network information. Their disclosure supports market transparency, coordination and public oversight, while some quantities remain restricted or become public only after a delay. For example, electricity-market disclosure rules distinguish forecasts and requirements before clearing, prices after clearing, and actual operating quantities released on the following day~\cite{nea2024disclosure}. Data classification standards likewise require field-level grades to account for aggregation and processing~\cite{gbt43697}. The practical question is how to give a field a meaningful grade when its information content depends on the other fields with which it can be combined.

This dependence follows from the system that generates the data. A generating unit's output identifies its operating point; a price reflects the marginal offer and network conditions. Either series may reveal little about a private offer parameter, yet together they can expose it. Congestion can weaken this relation for a remote price while preserving it at the unit's own bus. The controlled case in Section~\ref{sec:res-mech} makes this mechanism explicit. More generally, the background held by an adversary can turn an apparently innocuous release into useful information about unit outputs, line flows or actual load.

The underlying inference and aggregation problem is well established in database security~\cite{lunt1989aggregation,farkas2002inference}. It is also relevant to operational power data: public locational prices can reveal non-public network information~\cite{kekatos2014topology}. Recent inference-graph methods recognize that sensitivity depends on relations between fields~\cite{hu2026ignn}. A remaining difficulty is to connect a field's grade directly to what can be inferred from combinations of data, including combinations whose members have little individual predictive value.

We address this difficulty by starting with set-level inferability: the out-of-sample accuracy of an attacker retrained on a specified field set. Adding a field to different backgrounds produces a distribution of accuracy gains. Our audit retains the largest gain within a background budget, together with the single-field gain and the backgrounds that explain the difference. This is a budget-constrained form of the marginal contribution index (MCI)~\cite{catav2021mci}, evaluated relative to information already public. Its excess over single-field risk is the \emph{background amplification}. The resulting profile answers both how much a field can add and when that addition occurs.

Evaluating this profile requires many related inference tasks. With 35 candidates and at most two background fields, the audit needs 7,175 field sets per target. Dedicated retraining provides the reference evidence but makes repeated audits expensive. We therefore learn a shared representation once and fit a closed-form readout for each field set. Reconstruction-trained features and features from an untrained network provide complementary inputs to this readout. The estimator scans the candidate backgrounds; retraining the most promising ones verifies the reported risk evidence. Fig.~\ref{fig:framework} summarizes this progression from operational data to field profiles and release review.

The paper makes three contributions.
\begin{itemize}
\item A background-aware field profile derived from set-level inferability. We establish the safety-margin interpretation of budgeted marginal risk, characterize amplification with one background field, and connect critical fields to small minimal unsafe combinations.
\item An efficient audit procedure that estimates the full profile and verifies selected backgrounds by retraining. We evaluate its fidelity, background recovery, decision quality and computational cost, including a comparison with a tabular foundation model.
\item Evidence linking the measurement to power-system behavior. A controlled dispatch case identifies a known physical mechanism; four data sets with rule-based field roles establish the prevalence of amplification across ten targets. A withholding application shows how the resulting evidence can support release review.
\end{itemize}

\begin{figure}[!tb]
\centering\includegraphics[width=\textwidth]{fig1_framework.pdf}
\caption{From operational data to a background-aware field profile. A candidate field is evaluated over alternative backgrounds, producing a distribution of contextual gains. The audit separates the single-field contribution from amplification and records high-gain backgrounds. A shared estimator scans the combinations; dedicated retraining verifies selected evidence for release review. The miniature profile is illustrative.}\label{fig:framework}
\end{figure}
